% !PRINTED 373
\weberpage{373}
Nun kann man aber $b$ so wählen, daß
\weq{17}{b^2 \equiv D \ (\mathrm{mod}\ 4a^n)}
ist, für ein beliebig großes $n$. Denn man kann $b$ um ein beliebiges Vielfaches von $2a$ verändern. Nehmen wir also (17) als erfüllt an für einen Exponenten $n$, und setzen demgemäß
\weq{18}{b^2 = D + 4a^n c,}
so folgt:
\[
(b + 2ha^n)^2 - D \equiv 4a^n(c+hb) \ (\mathrm{mod}\ 4a^{n+1}),
\]
und wenn man also $h$ aus der Kongruenz $c + hb \equiv 0 \ (\mathrm{mod}\ a)$ bestimmt, so ergibt sich ein der Kongruenz (17) für $n+1$ genügendes $b$.

Da nun $b$ relativ prim zu $a$ ist, so folgt nach 2., daß, wenn $n \geqq m$ ist, $a^m$ die kleinste durch $\mathfrak{a}^m$ teilbare natürliche Zahl ist, und daraus ergibt sich folgende Kompositionsregel:

5. \gesperrt{Genügt die Form $(a, b, a^n c) = U$ der Bedingung (18), und ist $m \leqq n$, so ist}
\[
U^m = (a^m, b, a^{n-m}c).
\]
Hiermit sind die Klassen quadratischer Formen der Diskriminante $\Delta$ zu einer \gesperrt{Abelschen Gruppe} vereinigt, die mit der Gruppe der \gesperrt{Idealklassen isomorph} ist, und es ist zugleich ein Weg angegeben, wie man durch passend gewählte Repräsentanten die Komposition in diesen Gruppen wirklich ausführen kann\footnotemark[1]).

\begin{center}
\S\,102. \gesperrt{Komposition der Ordnungen.}
\end{center}

Eine Ordnung $O = [Q]$ im Körper $\Omega$ ist durch den Führer $Q$ vollständig bestimmt. Es seien $O_1$, $O_2$ zwei Ordnungen mit den Führern $Q_1$, $Q_2$, und der größte gemeinschaftliche Teiler von $Q_1$ und $Q_2$ sei $Q$. Ist dann $M$ das kleinste gemeinschaftliche Multiplum von $Q_1$ und $Q_2$, so ist
\weq{1}{Q_1 Q_2 = QM.}

\begingroup
\renewcommand{\thefootnote}{1)}
\footnotetext[1]{Die \gesperrt{Einheit} in der Gruppe der Klassen bildet die \gesperrt{Hauptklasse}. Über die Komposition der quadratischen Formen ist zu erwähnen: \gesperrt{Gauss}, \glqq Disquisitiones arithmeticae\grqq{}, Art. 235 u.\,f. \gesperrt{Dirichlet}, De formarum binariarum secundi gradus compositione (1851), Werke Bd. II, S. 105. \gesperrt{Dirichlet-Dedekind}, Vorlesungen über Zahlentheorie, 4. Aufl., \S\,145\,f. \gesperrt{Dedekind}, Crelles Journal, Bd. 129. \gesperrt{H. Weber}, Göttinger Nachrichten, 9. Februar 1907.}
\endgroup
